We study our method's sensitivity to batch size in Fig.~\ref{fig:batch_size_sensitivity}, by plotting the precision and recall achieved after convergence for many random initializations at many possible settings of batch size (from 512 up to the entire dataset), while training on the MIMIC mortality risk prediction task in the main paper. From Fig.~\ref{fig:batch_size_sensitivity}, we conclude that while all batch sizes have some initializations outperform others in terms of precision and recall, overall the best runs using batch sizes of 512 (orange dots) and 1024 (green dots) seem to reach higher precision and recall values than the best runs at other batch sizes. Only runs with 512 and 1024 have precision exceed 0.25 on heldout data (validation and test). We suspect that the ``folk wisdom'' that small batches add helpful ``noise'' to stochastic gradient descent in order to avoid bad local optima may be the reason that 512 and 1024 work better than the larger tested batch sizes.